\documentclass[a4paper,11pt]{ltjsarticle}

% 日本語対応
\usepackage{luatexja}
\usepackage{geometry}
\geometry{margin=25mm}

% 表・色
\usepackage{array}
\usepackage{longtable}
\usepackage[table]{xcolor}
\usepackage{booktabs}

% 箇条書き
\usepackage{enumitem}

% リンク
\usepackage{hyperref}

% セクションデザイン
\usepackage{titlesec}
\titleformat{\section}{\large\bfseries}{}{0em}{}

\setlength{\parindent}{0pt}
\setlist[itemize]{leftmargin=2em}

\begin{document}

\begin{center}
    {\LARGE \textbf{作業ログ}}\\[1em]
\end{center}

\begin{tabular}{|p{4cm}|p{10cm}|}
\hline
報告者 & 手代木巧\\
\hline
報告日 & 5月26日\\
\hline
プロジェクト名 & 赤外線通信で歌うカエル\\
\hline
\end{tabular}

\vspace{1em}

%=====================================================
\section{1. 進捗サマリー（要約）}

\begin{itemize}
    \item 全体進捗率：0％\
    \item マイルストーン達成状況：０個
    \item 主要成果：
    \begin{itemize}
        \item なし
        \item \underline{\hspace{1cm}}
    \end{itemize}
    \item 重大課題（影響度：高）：
    \begin{itemize}
        \item 進捗が０なこと
    \end{itemize}
\end{itemize}

%=====================================================
\section{2. 作業ログ（詳細トラッキング）}

\renewcommand{\arraystretch}{1.5}

\begin{longtable}{|p{2.2cm}|p{1.8cm}|p{2cm}|p{2cm}|p{1.5cm}|p{3cm}|p{2.5cm}|}
\hline
\rowcolor[gray]{0.9}
作業項目 & 開始日 & 予定完了日 & 状態 & 進捗率 & 依存関係・ブロッカー & コメント \\
\hline
回路作成 & 5/27& 6/3& 未着手 / 進行中 / レビュー中 / 完了 & & & \\
\hline
LEDマトリクスによるアニメーション & 6/3& 6/17& 未着手 / 進行中 / レビュー中 / 完了 & & & \\
\hline
\end{longtable}

%=====================================================
\section{3. 課題管理}

\begin{longtable}{|p{2cm}|p{2.5cm}|p{1.5cm}|p{1.5cm}|p{3cm}|p{2cm}|p{2cm}|}
\hline
\rowcolor[gray]{0.9}
課題 & 発生原因 & 影響度 & 優先度 & 対策 & 期限 & 状態 \\
\hline
課題1 & & 高 / 中 / 低 & 高 / 中 / 低 & & & 未対応 / 対応中 / 解決 \\
\hline
\end{longtable}

\vspace{0.5em}
※影響度＝プロジェクトへのインパクト\\
※優先度＝対応の緊急性

%=====================================================
\section{4. AI利用状況と検証（再現性重視）}

\subsection*{4.1 利用内容}

\begin{itemize}
    \item 利用目的：探索 / 設計 / 実装 / レビュー 等
    \item 入力（プロンプト要旨）：
    \begin{itemize}
        \item 何をどのように聞いたか
    \end{itemize}
    \item 出力概要：
    \begin{itemize}
        \item 何が得られたか
    \end{itemize}
\end{itemize}

\subsection*{4.2 検証方法}

※第三者が再現できるレベルで記述

\begin{itemize}
    \item 検証手法：
    \begin{itemize}
        \item 実験 / 比較 / 理論確認 など
    \end{itemize}
    \item 参照資料：
    \begin{itemize}
        \item 論文・書籍・公式ドキュメント等
    \end{itemize}
    \item 検証手順：
    \begin{enumerate}
        \item 手順1
        \item 手順2
    \end{enumerate}
\end{itemize}

\subsection*{4.3 ファクトチェック結果}

\begin{itemize}
    \item 一致点：
    \begin{itemize}
        \item 既存知識と整合している部分
    \end{itemize}
    \item 不一致・疑義：
    \begin{itemize}
        \item 差異の内容
    \end{itemize}
    \item 最終判断：採用 / 保留 / 棄却
    \item 根拠：
    \begin{itemize}
        \item 資料・理由
    \end{itemize}
\end{itemize}

%=====================================================
\section{5. 次回計画}

\begin{itemize}
    \item 目標：
    \begin{itemize}
        \item タスク内容
    \end{itemize}
    \item 達成基準：
    \begin{itemize}
        \item 検証可能な条件
    \end{itemize}
    \item 期限：
    \begin{itemize}
        \item 日付
    \end{itemize}
\end{itemize}

%=====================================================
\section{6. その他}

\begin{itemize}
    \item 連絡事項・懸念点など
\end{itemize}

%=====================================================
\section{7. 付録}

\begin{itemize}
    \item 添付資料：
    \begin{itemize}
        \item ソースコード：ファイル or リポジトリURL
        \item 回路図 / 設計図：ファイル
    \end{itemize}
    \item 添付範囲：
    \begin{itemize}
        \item 全体 / 差分（どちらか明記）
    \end{itemize}
\end{itemize}

\end{document}